\documentclass{beamer}

\usepackage{amsmath, amssymb, braket, tikz}
\usetikzlibrary{quantikz}

\title{Quantum Fourier Transform (QFT) on Two Qubits}
\author{MHJ}
\date{\today}

\begin{document}

\begin{frame}
   \titlepage
\end{frame}

% Slide 1: Introduction
\begin{frame}{Introduction}
   \textbf{Quantum Fourier Transform (QFT)}
   \begin{itemize}
       \item Quantum analogue of the Discrete Fourier Transform (DFT).
       \item Used in quantum algorithms like Shor's algorithm and phase estimation.
       \item Converts computational basis states into a superposition of phases.
   \end{itemize}
\end{frame}

% Slide 2: Definition
\begin{frame}{Definition of QFT}
   The QFT on an \( n \)-qubit system is defined as:
   \[
   \text{QFT} \ket{x} = \frac{1}{\sqrt{2^n}} \sum_{y=0}^{2^n-1} e^{2\pi i xy / 2^n} \ket{y}
   \]
   where \( x \) is the input integer in binary representation.
\end{frame}

% Slide 3: QFT on Two Qubits
\begin{frame}{QFT on Two Qubits}
   For a two-qubit system (\( n=2 \)), the QFT transformation is:
   \[
   \text{QFT} \ket{x_1 x_0} = \frac{1}{2} \sum_{y=0}^{3} e^{2\pi i x y / 4} \ket{y}
   \]
   where \( x = x_1 x_0 \) (binary representation).
\end{frame}



% Slide 5: Matrix Representation
\begin{frame}{Matrix Representation}
   The QFT on two qubits has the matrix:
   \[
   U_{\text{QFT}} = \frac{1}{2}
   \begin{bmatrix}
       1 & 1 & 1 & 1 \\
       1 & i & -1 & -i \\
       1 & -1 & 1 & -1 \\
       1 & -i & -1 & i
   \end{bmatrix}
   \]
\end{frame}

% Slide 6: Example Calculation
\begin{frame}{Example Calculation}
   Applying QFT to \( \ket{10} \) (binary 2):
   \[
   \text{QFT} \ket{10} = \frac{1}{2} (\ket{0} + i\ket{1} - \ket{2} - i\ket{3})
   \]
\end{frame}

% Slide 7: Conclusion
\begin{frame}{Conclusion}
   \begin{itemize}
       \item The QFT transforms basis states into Fourier superpositions.
       \item Two-qubit QFT involves Hadamards, phase gates, and a SWAP.
       \item It is essential for quantum algorithms like phase estimation.
   \end{itemize}
\end{frame}

\end{document}
% Slide 4: Quantum Circuit
\begin{frame}{Quantum Circuit for QFT (2 Qubits)}
   \begin{center}
       \begin{quantikz}
           \lstick{\ket{x_1}} & \gate{H} & \ctrl{1} & \qw      & \swap{1} & \qw \\
           \lstick{\ket{x_0}} & \gate{H} & \gate{R_2} & \qw & \targX{} & \qw
       \end{quantikz}
   \end{center}
   \begin{itemize}
       \item Hadamard $H$ applies equal superposition.
%       \item $R_2$ is the controlled phase gate: $R_2 = \begin{bmatrix} 1 & 0 \\ 0 & e^{i\pi/2} \end{bmatrix}$.
%       \item The final SWAP reorders the qubits.
   \end{itemize}
\end{frame}
